\section{Future Work}\label{sec:future}
Four directions follow directly from the findings. \emph{(i) A third source.} DA-GAFS selects with two labelled collections and is evaluated on held-out domains of the same collections; a confirmatory test must select on two sources and evaluate on a third, untouched one, under a new protocol version and holdout. \emph{(ii) Conditional direction.} The marginal direction gate removed a reversed but useful cue; replacing it with a conditional criterion, such as sign agreement of SHAP attributions within comparable slices, or penalising rather than excluding conflicts, is the natural refinement of C2. \emph{(iii) Contemporary and temporal data.} Both datasets are older public collections in which HTTPS is a strong benign cue; repeating the protocol on recent feeds with a temporal split would test whether this cue still transfers. \emph{(iv) Threshold transfer.} Ranking can partly transfer while the source threshold does not; label-free recalibration from target score distributions, compared with UDA~\cite{rashid2024} under the same domain-disjoint splits, would separate representation shift from operating-point shift. None of these may reuse the present TEST data for model choice.
